\documentclass{article}
\usepackage[T1]{fontenc}
\usepackage{lmodern}
\usepackage{graphicx}
\usepackage{amsmath,amssymb}
\usepackage[a4paper,margin=2cm]{geometry}
\usepackage[absolute,overlay]{textpos}
\setlength{\TPHorizModule}{1mm}
\setlength{\TPVertModule}{1mm}
\usepackage[indonesian]{babel}
\usepackage{hyperref}
\setlength{\emergencystretch}{2em}
% unit: Halaman 32

\begin{document}

% === Halaman 32 (python/native) ===

IP

IP-1 

Ki 

Si mod 8 

Linear Transformation

(except last round)

K32 

For i = 0 to 31 

32 copies of S-Box used.

4 bit input to each.

Linear 

Transformation

Output bits = 

 of input bits

Bit j = 0 to 127:

Odd j: XOR of 3 bits

Even j: XOR of 7 bits

Plaintext

Ciphertext

128 bit data block

32 rounds

256 bit keys, 

pads shorter keys 

Decryption differs from encryption

whitening

whitening

Sumber: Introduction to Practical Cryptography, Lecture 3 Block Ciphers 

32

\end{document}
